\begin{frame}[fragile]{JavaScript: Single Thread Model}
% https://dev.to/bbarbour/if-javascript-is-single-threaded-how-is-it-asynchronous-56gd
\begin{itemize}
	\item JavaScript uses `Single Thread Model.'
	\item However, in JavaScript, we can implement an asynchronous program that executes multiple functions at the same time.
	\item In this example, setTimeout() function executes the lambda expression in 1000 ms, but line 5 is executed without waiting for its results. 
\end{itemize}

\begin{columns}%
\begin{column}{0.4\textwidth}%
\begin{Code}{}%
\begin{lstlisting}[firstnumber=1, label=glabels, xleftmargin=0pt, basicstyle=\scriptsize]% 

console.log("first")
setTimeout(() => {
    console.log("second")
}, 1000)
console.log("third")
\end{lstlisting}%   
\end{Code}%
\end{column}%
\begin{column}{0.4\textwidth}%
\begin{Command}{Results}%
\begin{lstlisting}[firstnumber=1, label=glabels, xleftmargin=0pt, basicstyle=\tiny]% 

first
third
second
\end{lstlisting}%   
\end{Command}%%
\end{column}%
\end{columns}%
\end{frame}

\begin{frame}[fragile]{}

\begin{itemize}
	\item Understanding how JavaScript's single thread model can implement other programming languages' multi thread features is important. 
	\item However, there are many misunderstandings on JavaScript thread model that leads to confusion. 
\end{itemize}

\begin{center}
\includegraphics[width=0.5\textwidth]{thread.jpg}
\end{center} 
\end{frame}

\begin{frame}[fragile]{JavaScript Call Stack}
% https://www.javascripttutorial.net/javascript-event-loop/
\begin{itemize}
	\item The secret behind JavaScript single-thread model is a call stack and event loop. 
	\item The JavaScript engine executes a script from the top of the file and works its way down. 
	\item It creates the execution contexts, pushes, and pops functions onto and off the call stack in the execution phase.
\end{itemize}

\begin{center}
\includegraphics[width=0.4\textwidth]{eventloop1.png}
\end{center} 
\end{frame}

\begin{frame}[fragile]{}

\begin{itemize}
	\item In this example, console.log() function is pushed in the call stack and popped out (\circled{2}). 
	\item For time consuming function (task), it blocks other functions to be executed (\circled{4}), as console.log() function can be executed only after the task() function (\circled{6}). 
\end{itemize}

\begin{columns}%
\begin{column}{0.5\textwidth}%
\begin{Code}{}%
\begin{lstlisting}[firstnumber=1, label=glabels, xleftmargin=0pt, basicstyle=\tiny]% 

function task(message) {
    // emulate time consuming task
    let n = 10000000;
    while (n > 0){
        n--;
    }
    console.log(message);
}

console.log('Start script...');
task('Call an API');
console.log('Done!');
\end{lstlisting}%   
\end{Code}%
\end{column}%
\begin{column}{0.5\textwidth}%
\begin{center}%
\includegraphics[width=0.9\textwidth]{callstack.png} 
\end{center}%
\end{column}%
\end{columns}%
\end{frame}

\begin{frame}[fragile]{}

\begin{itemize}
	\item To address this block issue, we can use callback queue.
	\item In this example, when calling the setTimeout() function, the JavaScript engine places it on the call stack, and the Web API creates a timer that expires in 1 second in the Thread Pool.
\end{itemize}

\begin{columns}%
\begin{column}{0.5\textwidth}%

\begin{Code}{}%
\begin{lstlisting}[firstnumber=1, label=glabels, xleftmargin=0pt, basicstyle=\scriptsize]% 

console.log('Start script...');
setTimeout(() => {
    task('Download a file.');
}, 1000);
console.log('Done!');
\end{lstlisting}%   
\end{Code}%
\end{column}%
\begin{column}{0.5\textwidth}%
\begin{center}%
\includegraphics[width=0.9\textwidth]{javascript-event-loop-step-1.png}
\end{center}%
\end{column}%
\end{columns}%
\end{frame}

\begin{frame}[fragile]{Callback Queue and Event Loop}

\begin{itemize}
\item Then the JavaScript engine places the task() function into a callback queue. 
\item The event loop is a constantly running process that monitors both the callback queue and the call stack. 
\item If the call stack is not empty, the event loop waits until it is empty and places the next function from the callback queue to the call stack. 
\end{itemize} 

\begin{columns}%
\begin{column}{0.4\textwidth}%
\begin{center}%
\includegraphics[width=0.5\textwidth]{javascript-event-loop-step-2.png}
\end{center}%
\end{column}%
\begin{column}{0.4\textwidth}%
\begin{center}%
\includegraphics[width=0.5\textwidth]{javascript-event-loop-step-3.png} 
\end{center}%
\end{column}%
\end{columns}%
\end{frame}

\begin{frame}[fragile]{}

\begin{itemize}
\item When setTimeout() function is expired immediately, the 2nd console.log() function is placed in the call stack.
\item As a result, the task() in the setTimeout() function is executed the last. 
\end{itemize}

\begin{columns}%
\begin{column}{0.7\textwidth}%
\begin{center}%
\begin{Code}{}%
\begin{lstlisting}[firstnumber=1, label=glabels, xleftmargin=0pt, basicstyle=\scriptsize]%  

console.log('Hi!');

setTimeout(() => {
    console.log('Execute immediately.');
}, 0);

console.log('Bye!');
\end{lstlisting}%   
\end{Code}%
\end{center}%
\end{column}%
\begin{column}{0.3\textwidth}%
\begin{Command}{Results}%
\begin{lstlisting}[firstnumber=1, label=glabels, xleftmargin=0pt, basicstyle=\tiny]% 

Hi!
Bye!
Execute immediately.
\end{lstlisting}%   
\end{Command}%
\end{column}%
\end{columns}%

\end{frame}